\documentclass[conference]{IEEEtran}

\usepackage[utf8]{inputenc}
\usepackage[T1]{fontenc}
\usepackage{lmodern}
\usepackage{amsmath}
\usepackage{amssymb}
\usepackage{booktabs}
\usepackage{array}
\usepackage{tabularx}
\usepackage{graphicx}
\usepackage{hyperref}
\usepackage{xcolor}
\usepackage{listings}

\hypersetup{
    colorlinks=true,
    linkcolor=black,
    urlcolor=blue,
    citecolor=black
}

\lstdefinestyle{pythonstyle}{
    basicstyle=\ttfamily\footnotesize,
    breaklines=true,
    frame=single,
    columns=fullflexible
}

\title{Design and Implementation of a Desktop Analyzer for FNIRSI 1014D Oscilloscope Waveform Exports}

\author{
\IEEEauthorblockN{Project Technical Report}
\IEEEauthorblockA{
Software documentation prepared for the repository\\
\texttt{osciloscope\_fnirsi\_1014d\_analizer}
}
}

\begin{document}

\maketitle

\begin{abstract}
This document describes the design, implementation, and current capabilities of a signal analysis application for waveform captures exported by the FNIRSI 1014D oscilloscope. The system reads binary \texttt{.wav} files generated by the instrument, reconstructs both acquisition channels in physical voltage units, and exposes a set of advanced analysis modules through a desktop-oriented graphical interface built with Flask and PyWebView. The software supports oscilloscope configuration decoding, automatic measurement reconstruction, mathematical operations between channels, fast Fourier transform based spectral analysis, statistical analysis, advanced temporal measurements, numerical derivative and integral computation, correlation analysis, manual cursor measurements, per-cycle analysis, snapshots, LaTeX export, PNG export, and Windows desktop packaging through PyInstaller. In addition to the implementation details, this report presents a validation methodology, a reproducibility framework, suggested experimental evaluation criteria, and a discussion of the main technical limitations required for a stronger IEEE-style engineering report.
\end{abstract}

\begin{IEEEkeywords}
oscilloscope, waveform analysis, FFT, signal processing, Flask, PyWebView, PyInstaller, FNIRSI 1014D
\end{IEEEkeywords}

\section{Introduction}
Low-cost digital oscilloscopes are increasingly used in laboratory, academic, and field-debugging environments. However, the analysis features provided by embedded instrument firmware are often limited to a small subset of measurements and fixed visualization workflows. When a device such as the FNIRSI 1014D exports captured data to a \texttt{.wav} file, a software layer can extend the utility of that data by reconstructing signals offline and providing additional analysis modules that are either unavailable or inconvenient on the instrument itself.

The project documented in this report addresses that need. It implements a desktop application capable of parsing FNIRSI 1014D waveform exports, restoring the captured channels, and offering a professional analysis environment focused on rapid inspection, laboratory reporting, and engineering interpretation. The system is structured as a Python application with a Flask backend, a Jinja/HTML frontend, numerical processing based on NumPy and SciPy, visualization through Matplotlib, and a desktop wrapper powered by PyWebView. The resulting application can be executed either as a local Python program or as a packaged Windows executable.

\section{Project Objectives}
The project was designed around the following technical objectives:

\begin{enumerate}
    \item Decode FNIRSI 1014D waveform exports and reconstruct valid time-domain signals for both oscilloscope channels.
    \item Recover oscilloscope configuration metadata and built-in automatic measurements from the binary file format.
    \item Provide engineering analysis modules beyond the original instrument capabilities.
    \item Export results in forms appropriate for academic and technical reporting, particularly PNG figures and LaTeX fragments.
    \item Offer a desktop execution model suitable for users who do not want to operate the software from a command-line-only environment.
\end{enumerate}

\section{Related Work and Technical Context}
Offline analysis of instrument exports sits at the intersection of digital signal processing, laboratory software, and lightweight engineering tooling. From a processing perspective, the project relies on standard DSP concepts such as discrete-time Fourier analysis, cross-correlation, finite-difference derivative estimation, and numerical integration, all of which are well established in the literature \cite{oppenheim2010,proakis2007,lyons2011}. 

From a software perspective, the project follows the pattern of scientific computing environments built on NumPy, SciPy, and Matplotlib \cite{numpy2020,scipy2020,hunter2007}. This stack is especially appropriate for engineering prototypes because it combines numerical performance, transparent formulas, and mature plotting facilities.

The specific niche addressed by this application is narrower: converting waveform exports from a low-cost oscilloscope into a richer analysis environment. Commercial laboratory tools and vendor-specific applications often provide export viewing and measurement support, but they are frequently closed, limited to the features exposed by the original instrument, or poorly adapted to flexible reporting workflows. The relevance of the present project lies in bridging that gap for the FNIRSI 1014D by combining binary parsing, extended numerical analysis, export-oriented reporting, and desktop packaging in a single system.

\section{System Overview}
At a high level, the application follows the workflow shown below:

\begin{enumerate}
    \item A user loads a \texttt{.wav} capture exported by the FNIRSI 1014D.
    \item The backend validates the file and extracts oscilloscope configuration, automatic measurements, and raw channel samples.
    \item Raw samples are converted to voltage-domain signals using the decoded scope settings.
    \item Optional preprocessing, calibration, and mathematical combinations are applied.
    \item The user activates analysis modules on demand.
    \item The system renders graphs, numerical summaries, and exportable LaTeX/PNG outputs.
\end{enumerate}

The current implementation deliberately evaluates advanced modules only when the user requests them. This reduces unnecessary computation and improves interactivity, especially for FFT, statistics, advanced measurements, correlation, derivative/integral analysis, and cycle analysis.

\begin{figure}[ht]
\centering
\fbox{\parbox{0.92\columnwidth}{
\centering
\vspace{0.8em}
\textbf{Figure Placeholder 1. System architecture and processing flow}\\
\vspace{0.4em}
Recommended final content: upload stage, parser layer, signal reconstruction, analysis modules, visualization layer, LaTeX/PNG export, and desktop packaging path.
\vspace{0.8em}
}}
\caption{Suggested architecture figure to be replaced with a final diagram or workflow image.}
\label{fig:architecture-placeholder}
\end{figure}

\section{Repository Structure}
The main repository components are summarized in Table~\ref{tab:repo-structure}.

\begin{table}[ht]
\caption{Main Repository Files and Responsibilities}
\label{tab:repo-structure}
\centering
\begin{tabularx}{\columnwidth}{>{\raggedright\arraybackslash}p{0.28\columnwidth} X}
\toprule
\textbf{File} & \textbf{Responsibility} \\
\midrule
\texttt{app.py} & Main Flask application, session state, upload handling, partial updates, desktop window launch, and export endpoints. \\
\texttt{file\_analizer.py} & Binary parsing of FNIRSI files, configuration extraction, automatic measurements, and raw display data loading. \\
\texttt{signal\_analyzer.py} & Numerical processing pipeline including conversion, filtering, FFT, statistics, advanced metrics, correlation, cursors, and calibration. \\
\texttt{plot\_maker.py} & Time-domain, FFT, correlation, and auxiliary graph generation with Matplotlib. \\
\texttt{report.py} & LaTeX fragment generation and wrappers for downloadable PNG outputs. \\
\texttt{templates/main.html} & Main user interface, controls, module panels, toasts, cursor interaction, and desktop download integration. \\
\texttt{static/style.css} & Application styling and visual system. \\
\texttt{app.spec} & PyInstaller specification for packaging the application as a Windows executable. \\
\texttt{requirements.txt} & Runtime dependency declaration. \\
\bottomrule
\end{tabularx}
\end{table}

\section{Software Architecture}
\subsection{Execution Model}
The application uses Flask as the application server and PyWebView as the desktop container. In desktop mode, the backend is started in a local thread on \texttt{127.0.0.1:5000}, after which a native window is opened and pointed to the local address. This approach preserves the simplicity of a web-based UI while packaging the system as a desktop application.

The packaging configuration in \texttt{app.spec} includes the template and static asset folders so that the project can be distributed as \texttt{dist/app.exe} under Windows.

\subsection{State Management}
The project stores active analysis state in the Flask session and keeps numerical and plotting caches in memory to reduce repeated recomputation. This is particularly important because the user interface includes multiple advanced modules and high-resolution plots.

The application also manages temporary uploads explicitly:

\begin{itemize}
    \item uploaded captures are stored in \texttt{uploads/},
    \item reset operations clean the session and remove uploaded files,
    \item desktop window close events also trigger upload cleanup,
    \item an \texttt{atexit} safeguard performs additional cleanup on normal process exit.
\end{itemize}

\subsection{Partial UI Updates}
The final version of the project includes AJAX-oriented partial refreshes for analysis modules. Instead of rebuilding the visible page through a classic full-page postback every time a module action is triggered, the backend returns targeted fragments for the active panel. This design reduces perceived interface jumps and keeps the visual experience more stable while preserving the server-rendered architecture.

\subsection{Design Rationale}
Several implementation decisions are directly tied to usability and robustness goals:

\begin{itemize}
    \item \textbf{Flask + PyWebView:} This combination preserves a simple web technology stack while delivering a desktop experience suitable for non-web users.
    \item \textbf{On-demand analysis:} Expensive modules such as FFT and correlation are executed only when requested, reducing latency and avoiding unnecessary work.
    \item \textbf{In-memory caches:} Cached intermediate results reduce repeated parsing, reprocessing, and plot generation.
    \item \textbf{Upload cleanup:} Temporary file deletion on reset and on application close reduces stale-state accumulation and avoids user confusion.
    \item \textbf{Numerical sanitization:} Non-finite values are cleaned before advanced processing so that chained analysis remains stable.
\end{itemize}

Relative to alternatives such as full client-side rendering or a purely local GUI toolkit, this architecture offers a favorable trade-off between development speed, modularity, portability, and engineering transparency.

\section{Waveform File Decoding}
\subsection{Binary Parsing Strategy}
The module \texttt{file\_analizer.py} reads the FNIRSI waveform export as a binary structure. It extracts:

\begin{itemize}
    \item voltage-per-division configuration for both channels,
    \item probe attenuation,
    \item coupling mode,
    \item time-per-division,
    \item automatic oscilloscope measurements,
    \item raw display waveform data.
\end{itemize}

To improve robustness, the parser uses fixed-size reads and index validation. The custom \texttt{ScopeFileError} exception is raised when the file is truncated or contains out-of-range configuration indexes. This avoids silent corruption and makes invalid-file handling explicit at the application layer.

\subsection{Recovered Measurements}
The parser reconstructs the instrument measurements:

\begin{itemize}
    \item \texttt{Vmax}, \texttt{Vmin}, \texttt{Vavg}, \texttt{Vrms}, \texttt{Vpp}, \texttt{Vp},
    \item \texttt{Freq}, \texttt{Cycle},
    \item \texttt{Time+}, \texttt{Time-},
    \item \texttt{Duty+}, \texttt{Duty-}.
\end{itemize}

The current implementation further refines some timing values using duty-cycle information instead of assuming a fixed 50\% duty ratio, improving consistency for non-symmetric waveforms.

\section{Signal Reconstruction and Processing Pipeline}
\subsection{Sampling Frequency and Time Axis}
The sampling frequency is derived from the number of captured samples and the oscilloscope horizontal scale:

\begin{equation}
f_s = \frac{N}{T_{\text{screen}}}
\end{equation}

where $N$ is the sample count and $T_{\text{screen}}$ is the total horizontal time covered by the displayed divisions. The time axis is then generated as:

\begin{equation}
t[n] = \frac{n}{f_s}, \qquad n = 0,1,\dots,N-1
\end{equation}

\subsection{Voltage Conversion}
The module \texttt{signal\_analyzer.py} converts raw scope data to physical voltage values using the decoded vertical settings and measurement references. This stage is essential because the binary waveform samples are not directly meaningful without reconstruction into the electrical domain.

\subsection{Signal Conditioning}
The processing layer includes helper functions for:

\begin{itemize}
    \item replacing non-finite values,
    \item filtering out invalid divisions in mathematical operations,
    \item adaptive filtering through SciPy-based signal processing,
    \item user calibration through gain, offset, inversion, and optional normalization.
\end{itemize}

This combination allows the application to remain numerically stable when users activate chained operations such as calibration, math, FFT, derivative analysis, and cursor measurement.

\section{Validation Methodology}
For an IEEE-style engineering report, implementation description alone is not sufficient. A formal validation process is required to demonstrate that the reconstructed waveforms and derived metrics are trustworthy. The current project codebase already supports the necessary calculations; however, a complete benchmark dataset and controlled evaluation campaign are still pending. This section therefore defines the recommended validation protocol and records the current validation status honestly.

\subsection{Recommended Test Signals}
The project should be validated with at least the following signal families:

\begin{itemize}
    \item pure sinusoids with known amplitude and frequency,
    \item square waves with multiple duty cycles,
    \item triangular waves,
    \item pulse trains with controlled widths,
    \item noisy periodic signals,
    \item two-channel delayed versions of the same waveform for correlation testing,
    \item real captures exported by the FNIRSI 1014D under laboratory conditions.
\end{itemize}

\subsection{Evaluation Targets}
The following outputs should be compared against theory, reference software, or instrument readings:

\begin{itemize}
    \item frequency,
    \item RMS voltage,
    \item peak-to-peak voltage,
    \item cycle time and duty cycle,
    \item dominant FFT frequency,
    \item harmonic amplitudes and THD,
    \item correlation peak and delay,
    \item rise time and fall time,
    \item derivative peak and integral final value.
\end{itemize}

\subsection{Reference Sources for Comparison}
Three reference levels are appropriate:

\begin{enumerate}
    \item \textbf{Analytical reference:} signals generated from known closed-form expressions.
    \item \textbf{Instrument reference:} automatic measurements reported by the oscilloscope itself.
    \item \textbf{Software reference:} a trusted DSP environment such as MATLAB, SciPy scripts, or another laboratory analysis tool.
\end{enumerate}

\subsection{Error Metrics}
For each measurable quantity, the error should be reported using both absolute and relative forms:

\begin{equation}
e_{\mathrm{abs}} = |x_{\mathrm{estimated}} - x_{\mathrm{reference}}|
\end{equation}

\begin{equation}
e_{\mathrm{rel}}(\%) = \frac{|x_{\mathrm{estimated}} - x_{\mathrm{reference}}|}{|x_{\mathrm{reference}}|} \times 100
\end{equation}

For noisy or repeated trials, the mean error, standard deviation, and maximum deviation should also be reported.

\subsection{Current Validation Status}
At the time of this document, the project has undergone implementation-oriented verification and iterative functional correction, but it does \textit{not yet include} a formal benchmark suite with controlled numerical results stored in the repository. Therefore, the software can already be documented as a complete engineering implementation, but any strict performance claims should be presented as pending experimental validation rather than as fully certified measurement accuracy.

\section{Experimental Plan and Result Tables}
To support a stronger academic version of this report, Tables~\ref{tab:validation-cases} and \ref{tab:results-template} define a practical evaluation matrix that can be filled with future measurements.

\begin{table}[ht]
\caption{Recommended Validation Cases}
\label{tab:validation-cases}
\centering
\begin{tabularx}{\columnwidth}{>{\raggedright\arraybackslash}p{0.18\columnwidth} >{\raggedright\arraybackslash}p{0.20\columnwidth} X}
\toprule
\textbf{Case} & \textbf{Signal Type} & \textbf{Main Metrics to Verify} \\
\midrule
VC1 & Sine wave & Frequency, RMS, FFT dominant peak \\
VC2 & Square wave & Duty cycle, rise/fall time, harmonics, THD \\
VC3 & Triangle wave & RMS, derivative profile, cycle stability \\
VC4 & Pulse train & Time+, Time-, cursor and cycle accuracy \\
VC5 & Delayed dual-channel signal & Correlation peak and delay estimate \\
VC6 & Noisy real capture & Statistical robustness and FFT stability \\
\bottomrule
\end{tabularx}
\end{table}

\begin{table}[ht]
\caption{Result Table Template for Future Experimental Validation}
\label{tab:results-template}
\centering
\begin{tabularx}{\columnwidth}{>{\raggedright\arraybackslash}p{0.16\columnwidth} >{\raggedright\arraybackslash}p{0.18\columnwidth} >{\raggedright\arraybackslash}p{0.18\columnwidth} >{\raggedright\arraybackslash}p{0.18\columnwidth} X}
\toprule
\textbf{Case} & \textbf{Metric} & \textbf{Reference} & \textbf{Measured} & \textbf{Error} \\
\midrule
VC1 & Frequency & --- & --- & --- \\
VC1 & RMS & --- & --- & --- \\
VC2 & Duty cycle & --- & --- & --- \\
VC2 & THD & --- & --- & --- \\
VC5 & Delay & --- & --- & --- \\
\bottomrule
\end{tabularx}
\end{table}

\begin{figure}[ht]
\centering
\fbox{\parbox{0.92\columnwidth}{
\centering
\vspace{0.8em}
\textbf{Figure Placeholder 2. Representative interface and analysis outputs}\\
\vspace{0.4em}
Recommended final content: one screenshot of the main waveform view, one FFT example, and one cursor/correlation example.
\vspace{0.8em}
}}
\caption{Suggested visual evidence to strengthen the final academic version of the report.}
\label{fig:ui-placeholder}
\end{figure}

\section{Implemented Analysis Modules}
\subsection{Base Oscilloscope Measurements}
The main panel presents decoded scope configuration and automatic measurements together with a reconstructed time-domain graph for channels X and Y. These views serve as the reference context for every additional module.

\subsection{Mathematical Operations}
The application supports the following vector operations between channels:

\begin{equation}
X + Y,\quad X - Y,\quad X \times Y,\quad X \div Y
\end{equation}

This feature is implemented in \texttt{apply\_math\_operation}. For division, the code avoids unstable singular behavior by treating near-zero denominators carefully and sanitizing non-finite values before subsequent analysis modules use the result.

The math result is handled as a virtual signal and can be:

\begin{itemize}
    \item plotted,
    \item measured,
    \item analyzed spectrally,
    \item exported as PNG,
    \item exported as LaTeX.
\end{itemize}

\subsection{FFT Spectral Analysis}
The FFT module provides a one-sided magnitude spectrum based on NumPy's real FFT. Before transformation, the signal is centered and multiplied by a selectable window:

\begin{equation}
X_w[k] = \mathcal{F}\{(x[n] - \bar{x})w[n]\}
\end{equation}

Windowing is used to reduce spectral leakage in finite-length observations, which is a standard practice in discrete spectral estimation \cite{oppenheim2010,lyons2011}.

Supported windows are:

\begin{itemize}
    \item Rectangular,
    \item Hann,
    \item Hamming,
    \item Blackman.
\end{itemize}

The module computes:

\begin{itemize}
    \item magnitude spectrum,
    \item dominant frequency,
    \item dominant amplitude,
    \item top peaks,
    \item harmonics,
    \item total harmonic distortion (THD).
\end{itemize}

Dominant frequency estimation is reinforced by complementary time-domain methods elsewhere in the project, improving robustness compared to a naive single-bin FFT peak only.

\subsection{Statistical Analysis}
The statistics module computes signal descriptors directly from the voltage vectors:

\begin{itemize}
    \item mean,
    \item standard deviation,
    \item variance,
    \item median,
    \item minimum and maximum,
    \item range,
    \item RMS,
    \item peak-to-peak value.
\end{itemize}

This module is useful for quantifying signal stability, noise, offset, and overall variability. It supports X, Y, and the MATH trace when enabled.

\subsection{Advanced Signal Measurements}
The advanced module estimates dynamic waveform properties such as:

\begin{itemize}
    \item rise time,
    \item fall time,
    \item overshoot,
    \item undershoot,
    \item slew rate,
    \item crest factor.
\end{itemize}

Rise and fall times are based on 10\% and 90\% thresholds derived from robust percentile levels. This avoids reliance on a single min/max sample pair and produces more practical values for signals affected by small noise components.

\subsection{Derivative and Integral}
The numerical calculus module uses discrete approximations:

\begin{equation}
\frac{dx}{dt} \approx \nabla x
\end{equation}

and cumulative trapezoidal integration:

\begin{equation}
\int x(t)dt \approx \sum \frac{x[n] + x[n-1]}{2}\Delta t
\end{equation}

These approximations are consistent with standard numerical analysis practice for uniformly sampled discrete-time signals \cite{proakis2007,scipy2020}.

The results are displayed as additional graphs and exported independently. This feature is useful for transient-rate analysis and cumulative signal interpretation.

\subsection{Correlation Analysis}
Cross-correlation between X and Y is computed to estimate similarity and delay:

\begin{equation}
r_{xy}[\ell] = \sum_n x[n]y[n-\ell]
\end{equation}

Cross-correlation is a natural tool for estimating temporal alignment and similarity between discrete sequences \cite{oppenheim2010,lyons2011}.

The module reports:

\begin{itemize}
    \item the correlation curve,
    \item maximum correlation,
    \item estimated delay between channels.
\end{itemize}

This analysis is particularly relevant in synchronization, phase comparison, and system identification tasks.

\subsection{Manual Cursor Measurement}
The cursor module provides interactive measurement of two user-defined time locations over a selected signal. It returns:

\begin{itemize}
    \item $t_1$ and $t_2$,
    \item $V_1$ and $V_2$,
    \item $\Delta t$,
    \item $\Delta V$,
    \item estimated frequency from cursor spacing.
\end{itemize}

The final implementation uses direct browser-side cursor movement to improve visual alignment and reduce interaction lag relative to earlier overlay-based approaches.

\subsection{Per-Cycle Analysis}
The cycle module estimates cycle count, average period, average frequency, average peak-to-peak, and average RMS. This is useful for repetitive signals where a single global measurement may obscure cycle-to-cycle variation.

\subsection{Calibration and Normalization}
Calibration tools allow the user to:

\begin{itemize}
    \item scale channel amplitude,
    \item add or subtract offset,
    \item invert polarity,
    \item normalize the channels.
\end{itemize}

This module is valuable when exported captures need manual correction or when a normalized representation is useful before subsequent analysis.

\subsection{Snapshots and Comparison}
The snapshot system stores reduced summaries of previous analyses inside the session. A later capture can be compared against a stored snapshot, providing deltas for major measurements such as peak-to-peak voltage and frequency. This feature supports before-versus-after workflows during debugging and experimentation.

\section{Visualization and User Interface}
\subsection{Plotting Layer}
The plotting module centralizes the graph style and supports:

\begin{itemize}
    \item time-domain waveform plots,
    \item FFT plots,
    \item correlation plots,
    \item derivative and integral plots,
    \item cursor-related export plots.
\end{itemize}

The application uses different color palettes for dark-screen visualization and white-background export images. This design improves readability both on screen and in printed or embedded documents.

\subsection{Frontend Behavior}
The main interface is implemented in \texttt{templates/main.html} and styled through \texttt{static/style.css}. The UI includes:

\begin{itemize}
    \item an upload area for \texttt{.wav} captures,
    \item an oscilloscope configuration table,
    \item a base measurement table,
    \item a primary waveform graph,
    \item advanced analysis panels,
    \item copy-to-LaTeX and download controls,
    \item toast-based user feedback for successful operations.
\end{itemize}

These visual notifications replace blocking browser alerts and provide a more polished desktop experience.

\section{Discussion}
From an engineering standpoint, the most relevant contribution of the project is not a novel DSP algorithm but the integration of several signal-analysis techniques into a workflow adapted to a specific oscilloscope export format. That integration matters in practice because the FNIRSI export alone does not provide a rich environment for iterative analysis or selective report generation.

The project also illustrates an important trade-off: desktop usability was prioritized without abandoning the simplicity of a browser-based UI stack. Flask plus PyWebView is not the most minimal runtime, but it allows rapid feature extension, export handling, and modular interface design. The use of on-demand processing and cache layers is particularly important because some analysis modules, especially FFT and plot rendering, become noticeably heavier when executed repeatedly over the same uploaded data.

Another important design choice is explicit numerical sanitization. In a laboratory workflow, it is common to encounter imperfect data, clipped segments, or mathematically problematic operations such as channel division near zero. The decision to replace or mask non-finite values before advanced processing reduces crashes and preserves continuity of use, even if some edge cases still require careful interpretation by the user.

\section{Export and Reporting Features}
The project does not generate a monolithic report automatically. Instead, it offers section-oriented export suitable for flexible technical documentation:

\begin{itemize}
    \item PNG graph export for time-domain, MATH, FFT, derivative, integral, correlation, and cursor analysis,
    \item LaTeX fragment export for measurements, MATH, FFT, statistics, advanced measures, correlation, cursors, cycles, calibration, and comparisons.
\end{itemize}

This strategy is practical for laboratory reports because the user can selectively insert only the sections relevant to a given experiment.

\section{Desktop Packaging}
The system is intended for local desktop execution. In its packaged form:

\begin{itemize}
    \item Flask runs locally,
    \item PyWebView opens a native application window,
    \item PyInstaller bundles the application into a Windows executable,
    \item native save dialogs are used for downloads inside the desktop environment.
\end{itemize}

This solves a common issue in embedded browser contexts, where standard HTML file downloads may not behave like a full desktop browser.

\section{Current Strengths}
The project exhibits several strong engineering characteristics:

\begin{itemize}
    \item clear separation between parsing, analysis, plotting, export, and interface logic,
    \item broad analysis feature set built on top of oscilloscope exports,
    \item robust handling of invalid numerical values,
    \item desktop-oriented workflow with explicit upload cleanup,
    \item LaTeX and PNG export focused on real reporting needs,
    \item improved responsiveness through cache usage and partial updates.
\end{itemize}

\section{Known Limitations}
Despite the current maturity of the tool, several limitations remain:

\begin{itemize}
    \item the parser is tailored specifically to the FNIRSI 1014D export structure,
    \item the software depends on local execution and is not optimized for multi-user deployment,
    \item the project does not yet include a formal automated test suite,
    \item advanced timing and spectral metrics remain sensitive to signal quality, periodicity, and file correctness,
    \item the UI still mixes server-rendered content with client-side interactivity, which adds implementation complexity.
\end{itemize}

\subsection{Quantified Technical Constraints}
The following constraints should be made explicit in any formal evaluation:

\begin{itemize}
    \item frequency accuracy depends on sample count, time-window duration, signal-to-noise ratio, and the chosen FFT window;
    \item harmonic amplitude estimation and THD are sensitive to spectral leakage and to the exact alignment between waveform period and record length;
    \item rise/fall estimates are less reliable for strongly noisy or non-monotonic transitions;
    \item correlation delay estimation deteriorates when signals have weak similarity, severe distortion, or poor time alignment;
    \item the parser assumes that the vendor export structure remains stable across firmware versions;
    \item exported waveform resolution is inherently limited by the device capture and encoding process.
\end{itemize}

\section{Future Work}
The following extensions would further improve the project:

\begin{itemize}
    \item automated unit and regression tests for parsers and analysis functions,
    \item formal benchmarking of FFT and advanced measurement accuracy,
    \item additional export targets such as complete PDF-oriented pipelines,
    \item support for more oscilloscope file formats,
    \item deeper client-side updates to further reduce interface redraw cost.
\end{itemize}

\section{Conclusion}
This project demonstrates that waveform exports from a low-cost oscilloscope can be transformed into a significantly more capable analysis environment through a carefully structured Python application. By combining binary parsing, calibrated signal reconstruction, advanced numerical analysis, technical visualization, LaTeX export, and Windows desktop packaging, the software extends the practical usefulness of the FNIRSI 1014D beyond its native firmware interface.

The final result is not merely a waveform viewer; it is a desktop analysis workstation tailored to oscilloscope export data. Its architecture is sufficiently modular to support future research, refactoring, or feature growth, while already delivering a meaningful tool for laboratory work, signal interpretation, and technical reporting.

\section{Reproducibility and Software Validation Notes}
The present repository already contains the minimum information required to reproduce the software build and execution path. A recommended reproducibility checklist is the following:

\begin{enumerate}
    \item Use Python 3.11 in a Windows environment.
    \item Install the dependencies listed in \texttt{requirements.txt}.
    \item Execute \texttt{python app.py} for local desktop mode.
    \item Load an exported FNIRSI 1014D \texttt{.wav} capture.
    \item Run each analysis module on demand and store exported PNG and LaTeX outputs.
    \item Build the desktop executable through \texttt{pyinstaller app.spec}.
\end{enumerate}

For a publication-grade validation package, the repository should additionally include:

\begin{itemize}
    \item a curated set of test waveform files,
    \item expected numerical results,
    \item automated comparison scripts,
    \item versioned environment information,
    \item screenshots or generated figures referenced from the final paper.
\end{itemize}

\section*{Appendix: Main Runtime Dependencies}
\begin{lstlisting}[style=pythonstyle]
Flask
Werkzeug
numpy
scipy
matplotlib
pywebview
\end{lstlisting}

\section*{Appendix: Compilation Notes}
This document is intended to be compiled with the \texttt{IEEEtran} class. A typical compilation command in a LaTeX environment is shown below:

\begin{lstlisting}[style=pythonstyle]
pdflatex project_report_ieee.tex
\end{lstlisting}

\begin{thebibliography}{99}

\bibitem{oppenheim2010}
A. V. Oppenheim and R. W. Schafer, \textit{Discrete-Time Signal Processing}, 3rd ed. Upper Saddle River, NJ, USA: Pearson, 2010.

\bibitem{proakis2007}
J. G. Proakis and D. G. Manolakis, \textit{Digital Signal Processing: Principles, Algorithms, and Applications}, 4th ed. Upper Saddle River, NJ, USA: Pearson Prentice Hall, 2007.

\bibitem{lyons2011}
R. G. Lyons, \textit{Understanding Digital Signal Processing}, 3rd ed. Upper Saddle River, NJ, USA: Prentice Hall, 2011.

\bibitem{numpy2020}
C. R. Harris \textit{et al.}, ``Array programming with NumPy,'' \textit{Nature}, vol. 585, pp. 357--362, Sep. 2020.

\bibitem{scipy2020}
P. Virtanen \textit{et al.}, ``SciPy 1.0: Fundamental algorithms for scientific computing in Python,'' \textit{Nature Methods}, vol. 17, pp. 261--272, Mar. 2020.

\bibitem{hunter2007}
J. D. Hunter, ``Matplotlib: A 2D graphics environment,'' \textit{Computing in Science \& Engineering}, vol. 9, no. 3, pp. 90--95, 2007.

\bibitem{flaskdocs}
Pallets, ``Flask Documentation,'' [Online]. Available: \url{https://flask.palletsprojects.com/}. Accessed: Mar. 19, 2026.

\bibitem{pywebviewdocs}
PyWebView, ``PyWebView Documentation,'' [Online]. Available: \url{https://pywebview.flowrl.com/}. Accessed: Mar. 19, 2026.

\bibitem{pyinstallerdocs}
PyInstaller Development Team, ``PyInstaller Documentation,'' [Online]. Available: \url{https://pyinstaller.org/}. Accessed: Mar. 19, 2026.

\end{thebibliography}

\end{document}
